\section{What cognitive restructuring asks of the patient}\label{sec:demands}
Consider the standard thought-record exercise. The patient must (a) \emph{notice} an automatic thought while it is occurring, (b) hold the thought, the supporting evidence, the contradicting evidence and a candidate alternative in mind at once, (c) suppress the pull of the automatic thought while evaluating it, and (d) shift from being inside the feeling to observing it. Mapped onto the best-known model of executive function, these are three separable components: updating and holding information (working memory), inhibiting a prepotent response, and shifting between mental sets \citep{miyake2000unity}. Executive function is, to a first approximation, supported by prefrontal networks.

Two clarifications keep this claim honest. First, it is a task analysis, not an empirical result: the argument needs the demands to be \emph{real} in practice, and clinicians teach scaffolds (worksheets, a therapist who prompts the evidence, repetition) that lower the load. Second, executive function is not a single resource; the argument concerns the three components above, not a general ``intelligence'' required for therapy. We return to both points in Section~\ref{sec:objections}.
